\documentclass{situationist-preprint}

\title{Bodily Continuity as Tie-Breaker}
\msid{PID-01-b}
\msfamily{PID-01}
\msdiscipline{Personal Identity and the Self}
\mstype{Companion manuscript}
\mskeywords{personal identity; branching; bodily continuity; closest continuer; phase sortal; animalism}

\begin{document}
\maketitle

\begin{abstract}
PID-01 holds that a person persists by psychological continuity, and that bodily continuity breaks ties when psychological continuity branches. Critics of tie-breakers say that any such clause is ad hoc: it makes identity depend on a feature chosen only to rescue a criterion from counterexamples. We argue that a tie-breaker is required, that the alternatives to a bodily one fail, and that the bodily one is not ad hoc. On the reading defended here, the tie is broken by the continuation of the original organism's life, not by the quantity of matter retained. So read, the clause follows from the status of \enquote{person} as a phase sortal of the organism, and it yields the verdicts of PID-02, PID-03 and PID-14.
\end{abstract}

\section{Why a tie-breaker is needed}

The answer of PID-01 to the question what makes a person the same over time is psychological continuity, with bodily continuity as a tie-breaker in cases of branching. The psychological continuity in question is non-branching continuity through a cause of the right kind, including connections of character and value, where the cause is of the right kind only if the later stages are realized in a thing of the same kind as the earlier ones; and \enquote{person} is a phase sortal of the organism, true of it while it has the capacities for rational and self-conscious thought (PID-05-a, Definition~2.2).

\begin{definition}\label{d:branch}
A process \emph{branches} if at some time two or more later persons are each psychologically continuous with an earlier person, through a cause of the right kind. Each such later person is a \emph{branch}.
\end{definition}

A pure non-branching criterion says that a person survives as a later person if and only if they are psychologically continuous and the process does not branch. Applied without qualification, it gives an intolerable verdict in the Branch-Line case \citep[ch.~10]{Parfit1984}. A scanner records the state of a person's brain without harming it, and a replica is built on Mars. The process branches, so by the pure criterion the woman who steps out of the scanner on Earth, in the same body, with an unbroken stream of consciousness, is not the woman who stepped in. Nothing has happened to her except that a machine has read her, and something has happened elsewhere. A criterion on which reading a person ends her existence is not a criterion of personal identity. PID-02-a accordingly holds that she continues, for the days remaining to her, as the original (PID-02-a, Proposition~3.2). Some clause must say which branch, if any, the original is. That clause is the tie-breaker.

\section{The tie-breaker}

\begin{definition}\label{d:life}
A branch $y$ \emph{continues the life} of an earlier person $x$ if $y$ is realized in an organism whose life (PID-05-a, Definition~2.1) is the continuation of the life of the organism in which $x$ was realized.
\end{definition}

\begin{principle}[Tie-breaker]\label{p:tie}
In a branching process, the original person is identical with a branch if and only if that branch, and no other, continues her life.
\end{principle}

\begin{proposition}\label{p:verdicts}
Principle~\ref{p:tie} yields the following verdicts. (i) In Branch-Line the original is the person on Earth, not the replica. (ii) If the original is destroyed after the replica is activated, she is the person on Earth until her destruction, and the replica is a distinct person. (iii) In symmetric fission she is neither successor. (iv) In nondestructive uploading she is the person who remains in her body.
\end{proposition}

\begin{proof}
In (i), (ii) and (iv) the process branches and exactly one branch is realized in the organism whose life continues hers; a replica or an upload is realized in a newly made organism or in a machine, whose existence is caused by the original's but whose life, if it has one, is not the continuation of hers. In (iii) each hemisphere is housed in a new body and the original organism does not survive the operation, so no branch continues her life and she is neither.
\end{proof}

Verdict (ii) is the case added in the revision of PID-02-a (Theorem~3.3); verdict (iii) is PID-03-a, Lemma~3.1; verdict (iv) agrees with PID-14-a, Proposition~3.4.

\section{Against the alternatives}

There are three other ways to break ties.

\paragraph{Degree of psychological continuity.} The closest continuer theory \citep{Nozick1981} lets the branch more closely continuous with the original be the original. In Branch-Line the replica is, by stipulation, exactly similar, so the degrees are equal and the closest continuer theory selects neither, which returns us to the intolerable verdict. Where the degrees differ slightly, the theory makes identity turn on the difference between a memory retained and a memory lost, which is not the sort of fact that could make the difference between a person's survival and her death.

\paragraph{Degree of bodily continuity.} A tie-breaker by the quantity of body or brain retained has the second defect in another form. If one successor receives fifty-one per cent of the original's neurons and the other forty-nine, a quantitative tie-breaker makes the first the original. But nothing about the person, her thought or her life distinguishes the two successors in a way that could ground identity with one and not the other. This is why Principle~\ref{p:tie} looks to the continuation of a life, which is an all-or-nothing matter in the cases at issue, rather than to the amount of matter retained.

\paragraph{Convention.} One may hold that, when psychology branches, there is no fact about who the original is until we decide. A convention that settled the matter by stipulation would not explain why the Branch-Line verdict is not open to decision: no community could reasonably decide that the woman on Earth had ceased to exist.

\section{Why the clause is not ad hoc}

A clause is ad hoc if its only support is that it saves a theory from counterexamples. Principle~\ref{p:tie} has independent support from two features of the account.

First, \enquote{person} is a phase sortal of the organism. A person, in the ordinary case, is an organism during the phase in which it has the capacities for rational and self-conscious thought, and the person's stages are unified by psychological continuity because those capacities are exercised by one organism (PID-05-a, Proposition~4.1). When psychology branches, the organism is not thereby divided: it remains where its life continues. The person phase of that organism continues with it, so long as the organism retains the capacities. Principle~\ref{p:tie} says no more than this. It is the default from which the psychological criterion departs only where the organism ceases and its psychology continues uniquely elsewhere, as in Simple Teletransportation.

Second, the kind condition on the cause of continuity already gives the realizer a role in identity: continuity sustains the persistence of a person only when its later stages are realized in a thing of the same kind as the earlier ones. In a branching case each branch may meet this condition, and the question is which of two realizers of the right kind is the one the person had. The continuation of a single life is the natural answer, because a life is what makes successive stages of an organism stages of one realizer.

\begin{objection}
The tie-breaker makes identity extrinsic. Whether the woman on Earth is the original depends on whether a replica exists on Mars only through the fact that the process branches; but once branching occurs, whether the replica is the original depends on what happens to the body on Earth.
\end{objection}

\begin{reply}
The extrinsicness is real, and PID-02-a accepts it (Corollary~3.4). It is the price of any account on which identity is one-one and the relation that grounds it can hold one-many. What it does not imply is that what matters in survival is extrinsic. Whether a later being has what matters depends on the intrinsic character of its relation to the original (PID-03-a, Assumption~2.6), and so the replica in Branch-Line has much of what matters though it is not the original (PID-06).
\end{reply}

\begin{objection}
If bodily continuity can decide identity in branching cases, it should decide it everywhere, and the account collapses into a bodily criterion.
\end{objection}

\begin{reply}
The tie-breaker operates only among branches, each of which is psychologically continuous with the original through a cause of the right kind. It never makes a person identical with a body that has lost her psychology, nor does it deny survival where psychology continues uniquely without the original organism, as in teletransportation (PID-02-a, Proposition~3.1). A criterion on which bodily continuity is necessary would give the opposite verdict there.
\end{reply}

\section{Conclusion}

Without a tie-breaker the psychological criterion makes nondestructive scanning fatal. Ties cannot be broken by degrees of psychological or material continuity without making identity turn on trivial differences, nor by convention. The continuation of the original organism's life breaks them, gives the verdicts of PID-02, PID-03 and PID-14, and follows from the status of \enquote{person} as a phase sortal of the organism.

\begin{thebibliography}{99}\small
\bibitem[Nozick(1981)]{Nozick1981} Nozick, R. (1981). \emph{Philosophical Explanations}. Harvard University Press.
\bibitem[Parfit(1984)]{Parfit1984} Parfit, D. (1984). \emph{Reasons and Persons}. Oxford University Press.
\end{thebibliography}

\end{document}
